% Chapter 6 draft for textbook inclusion. Assumes a textbook preamble with amsmath, amssymb, and array.
\chapter{Simple Linear Regression: Fit Quality, $R^2$, and Hypothesis Testing}
\label{chap:slr-fit-quality-r2-hypothesis-testing}

\section{The Big Picture}

The previous chapter built the simple linear regression model and derived the least squares line
\begin{equation}
  \widehat{y}=\widehat{\beta}_0+\widehat{\beta}_1x.
\end{equation}
The next question is not how to find the line.  The next question is:
\begin{quote}
  Once we have a fitted line, how good is it, and is there statistical evidence that the predictor is useful?
\end{quote}

This chapter introduces two related but different ideas.
\begin{enumerate}
  \item \textbf{Fit quality:} How much variation in the response does the fitted line explain?  This is measured by $R^2$.
  \item \textbf{Statistical evidence:} Is the fitted relationship distinguishable from a zero-slope model after accounting for sample size and noise?  This is tested with the $F$-test or, equivalently in simple linear regression, the slope $t$-test.
\end{enumerate}

The difference matters.  A model can have a high $R^2$ and still be based on too few observations to trust.  A model can have a statistically significant slope but still explain little practical variation.  Good analysis reports both fit quality and inferential evidence.

\section{Review: The Mean-Only Baseline Model}

Before using a predictor, we can always estimate a quantitative response $Y$ by its sample mean:
\begin{equation}
  \widehat{y}_i=\bar{y},
  \qquad
  \bar{y}=\frac{1}{n}\sum_{i=1}^n y_i.
\end{equation}
This is the \emph{mean-only model}.  It ignores all predictor values.

The prediction error for observation $i$ under the mean-only model is
\begin{equation}
  y_i-\bar{y}.
\end{equation}
The mean squared error of the mean-only model is
\begin{equation}
  \frac{1}{n}\sum_{i=1}^n (y_i-\bar{y})^2.
\end{equation}
This is the denominator-$n$ version of sample variance.  Thus, for one-dimensional data, evaluating the mean-only estimate by MSE is essentially evaluating variance.

\paragraph{Math Foundations Connection.}
A response column $(y_1,\ldots,y_n)^T$ is a vector in $\mathbb{R}^n$, and the mean-only fitted vector can be written as $\bar{y}\mathbf{1}$.  The sums of squared deviations below are squared Euclidean norms; see the Math Foundations textbook, Chapters 2, 3, and 5.

\section{Adding a Predictor}

When we observe paired data
\begin{equation}
  (x_1,y_1),(x_2,y_2),\ldots,(x_n,y_n),
\end{equation}
we can compare two strategies:
\begin{enumerate}
  \item predict all responses using $\bar{y}$;
  \item predict responses using the fitted regression line
  \begin{equation}
    \widehat{y}_i=\widehat{\beta}_0+\widehat{\beta}_1x_i.
  \end{equation}
\end{enumerate}

If $X$ is useful, then the fitted line should have smaller squared prediction error than the mean-only model.  If $X$ is useless, then the fitted line will look like a horizontal line with slope near zero, and it will not improve much over the mean-only model.

This is the core intuition behind $R^2$ and the $F$-test:
\begin{quote}
  Compare the error from the mean-only model to the error from the fitted regression model.
\end{quote}

\section{Three Sums of Squares}

Let
\begin{equation}
  e_i=y_i-\widehat{y}_i
\end{equation}
be the residual from the fitted regression line.

\subsection{Total Sum of Squares}

The \emph{total sum of squares} measures how much the response varies around its sample mean:
\begin{equation}
  \operatorname{TSS}=\sum_{i=1}^n (y_i-\bar{y})^2.
  \label{eq:ch6-tss}
\end{equation}
This is the squared error from the mean-only model.

\subsection{Residual Sum of Squares}

The \emph{residual sum of squares} measures how much variation remains after fitting the regression line:
\begin{equation}
  \operatorname{RSS}=\sum_{i=1}^n (y_i-\widehat{y}_i)^2
  =\sum_{i=1}^n e_i^2.
  \label{eq:ch6-rss}
\end{equation}
This is the squared error from the fitted regression model.

\subsection{Regression Sum of Squares}

The \emph{regression sum of squares}, sometimes called explained sum of squares, measures how much variation the model explains relative to the mean-only baseline:
\begin{equation}
  \operatorname{RegSS}=\operatorname{TSS}-\operatorname{RSS}.
  \label{eq:ch6-regss}
\end{equation}

When the regression model includes an intercept, least squares gives the decomposition
\begin{equation}
  \boxed{\operatorname{TSS}=\operatorname{RegSS}+\operatorname{RSS}.}
  \label{eq:ch6-sum-of-squares-decomposition}
\end{equation}

\paragraph{Math Foundations Connection.}
The decomposition in~\eqref{eq:ch6-sum-of-squares-decomposition} is a geometry statement.  It relies on the residual vector being orthogonal to the fitted-value direction after centering.  Dot products, norms, and orthogonality are covered in the Math Foundations textbook, Chapters 2, 3, and 5.

\section{$R^2$: Proportion of Variation Explained}

The coefficient of determination is
\begin{equation}
  \boxed{R^2=\frac{\operatorname{TSS}-\operatorname{RSS}}{\operatorname{TSS}}
  =1-\frac{\operatorname{RSS}}{\operatorname{TSS}}.}
  \label{eq:ch6-r2}
\end{equation}

The numerator $\operatorname{TSS}-\operatorname{RSS}$ is the reduction in squared error achieved by using the predictor.  The denominator $\operatorname{TSS}$ is the squared error we had before using the predictor.

Thus,
\begin{quote}
  $R^2$ is the proportion of the response variation explained by the regression model.
\end{quote}

If $R^2=0$, then the fitted model does no better than the mean-only model.  If $R^2=1$, then the residuals are all zero and the fitted line passes exactly through every observed point.

For simple linear regression with an intercept,
\begin{equation}
  0\leq R^2\leq 1.
\end{equation}
This interval property depends on comparing an intercept-containing regression model to the mean-only baseline.  Some nonstandard models, such as no-intercept regressions, can produce different behavior.

\section{Interpreting $R^2$}

Suppose two models are used to predict pediatric weight.
\begin{itemize}
  \item Model A uses wrist circumference and has $R^2=0.25$.
  \item Model B uses height and has $R^2=0.50$.
\end{itemize}
Then Model B explains twice as much of the observed weight variation as Model A, at least in the data used to fit the models.

The phrase \emph{explains variation} should be read carefully.  It does not automatically imply causation.  It means that the fitted line reduces squared prediction error relative to the mean-only model.

\CoreCompress{
\section{$R^2$ and Correlation in Simple Linear Regression}

For simple linear regression with an intercept, $R^2$ is equal to the squared sample correlation between $x$ and $y$:
\begin{equation}
  \boxed{R^2=r_{xy}^2.}
\end{equation}

To see why, define centered values
\begin{equation}
  x_i^c=x_i-\bar{x},
  \qquad
  y_i^c=y_i-\bar{y}.
\end{equation}
Let
\begin{equation}
  S_{xx}=\sum_{i=1}^n (x_i^c)^2,
  \qquad
  S_{yy}=\sum_{i=1}^n (y_i^c)^2,
  \qquad
  S_{xy}=\sum_{i=1}^n x_i^cy_i^c.
\end{equation}
The least squares slope is
\begin{equation}
  \widehat{\beta}_1=\frac{S_{xy}}{S_{xx}}.
\end{equation}
The centered fitted values satisfy
\begin{equation}
  \widehat{y}_i-\bar{y}=\widehat{\beta}_1(x_i-\bar{x}).
\end{equation}
Therefore, the regression sum of squares is
\begin{align}
  \operatorname{RegSS}
  &=\sum_{i=1}^n (\widehat{y}_i-\bar{y})^2 \\
  &=\sum_{i=1}^n \widehat{\beta}_1^2(x_i-\bar{x})^2 \\
  &=\widehat{\beta}_1^2S_{xx} \\
  &=\left(\frac{S_{xy}}{S_{xx}}\right)^2S_{xx} \\
  &=\frac{S_{xy}^2}{S_{xx}}.
\end{align}
Thus,
\begin{align}
  R^2
  &=\frac{\operatorname{RegSS}}{\operatorname{TSS}} \\
  &=\frac{S_{xy}^2/S_{xx}}{S_{yy}} \\
  &=\frac{S_{xy}^2}{S_{xx}S_{yy}} \\
  &=\left(\frac{S_{xy}}{\sqrt{S_{xx}S_{yy}}}\right)^2 \\
  &=r_{xy}^2.
\end{align}

This equality is useful in SLR.  It does not directly generalize to multiple linear regression as a single ordinary correlation between one predictor and the response.
}{
\section{$R^2$ and Correlation in Simple Linear Regression}

For simple linear regression with an intercept,
\begin{equation}
  R^2=r_{xy}^2.
\end{equation}
This is a special one-predictor identity: the proportion of response variation explained by the fitted line equals the squared sample correlation between the predictor and the response.  In multiple regression, there is no single ordinary predictor-response correlation that plays the same role.
}
\section{Limitations of $R^2$}

$R^2$ is useful, but it does not tell us everything.

\subsection{A Perfect Fit Can Be Misleading}

With two data points and a line with two parameters, we can fit the data perfectly as long as the two $x$ values are different.  Then
\begin{equation}
  \operatorname{RSS}=0
  \qquad\Rightarrow\qquad
  R^2=1.
\end{equation}
This does not mean the model will predict future data well.  It means only that a line can pass through two points.

\subsection{$R^2$ Does Not Measure Statistical Evidence}

$R^2$ measures in-sample fit.  It does not directly answer:
\begin{itemize}
  \item Is the slope distinguishable from zero?
  \item How much sampling uncertainty is in $\widehat{\beta}_1$?
  \item How noisy are the observations?
  \item How many observations were used?
\end{itemize}

These questions require hypothesis tests, confidence intervals, and residual diagnostics.

\subsection{$R^2$ Does Not Check Model Assumptions}

A high $R^2$ does not prove that the linear model assumptions are valid.  Residual plots may still reveal curvature, nonconstant variance, outliers, or dependence among observations.

\section{From Fit Quality to Hypothesis Testing}

Fit quality asks:
\begin{quote}
  How much did the line reduce squared error compared to the mean-only model?
\end{quote}

Hypothesis testing asks:
\begin{quote}
  Is that reduction large enough, relative to residual noise and sample size, to count as statistical evidence that the slope is not zero?
\end{quote}

For SLR, the inferential question is usually
\begin{equation}
  H_0:\beta_1=0
  \qquad\text{versus}\qquad
  H_A:\beta_1\neq 0.
  \label{eq:ch6-slope-hypotheses}
\end{equation}
The null hypothesis says that $X$ does not improve the conditional mean of $Y$ through a linear slope.  The alternative says that the slope is nonzero.

\paragraph{Math Foundations Connection.}
Hypotheses are mathematical statements, and hypothesis testing depends on understanding logical alternatives and negation.  Basic logical statements, negation, and truth-table reasoning are covered in the Math Foundations textbook, Chapter 1.

\section{The $F$-Statistic}

The $F$-statistic compares the mean-only model to the fitted regression model.

For a general nested-model comparison,
\begin{equation}
  F=\frac{\left(\operatorname{SS}_{\text{null}}-\operatorname{SS}_{\text{full}}\right)/(p_{\text{full}}-p_{\text{null}})}
  {\operatorname{SS}_{\text{full}}/(n-p_{\text{full}})},
  \label{eq:ch6-general-f}
\end{equation}
where:
\begin{itemize}
  \item $\operatorname{SS}_{\text{null}}$ is the squared error from the simpler model;
  \item $\operatorname{SS}_{\text{full}}$ is the squared error from the larger model;
  \item $p_{\text{null}}$ is the number of parameters in the simpler model;
  \item $p_{\text{full}}$ is the number of parameters in the larger model.
\end{itemize}

For simple linear regression, the null model is the mean-only model, so
\begin{equation}
  p_{\text{null}}=1,
  \qquad
  \operatorname{SS}_{\text{null}}=\operatorname{TSS}.
\end{equation}
The full model is the fitted line, so
\begin{equation}
  p_{\text{full}}=2,
  \qquad
  \operatorname{SS}_{\text{full}}=\operatorname{RSS}.
\end{equation}
Therefore,
\begin{align}
  F
  &=\frac{(\operatorname{TSS}-\operatorname{RSS})/(2-1)}{\operatorname{RSS}/(n-2)} \\
  &=\boxed{\frac{\operatorname{TSS}-\operatorname{RSS}}{\operatorname{RSS}/(n-2)}}.
  \label{eq:ch6-slr-f}
\end{align}

The numerator is the improvement over the mean-only model.  The denominator is the residual variation per leftover degree of freedom.

\paragraph{Math Foundations Connection.}
The phrase \emph{leftover degrees of freedom} is a dimension idea.  The residual variation lives in directions left over after fitting the intercept and slope.  Linear independence, bases, and the fact that a linearly independent set of $n$-vectors can have at most $n$ elements are introduced in the Math Foundations textbook, Chapters 2, 5, and 6.

\section{The $F$-Statistic in Terms of $R^2$}

Because
\begin{equation}
  R^2=\frac{\operatorname{TSS}-\operatorname{RSS}}{\operatorname{TSS}},
\end{equation}
we can rewrite the SLR $F$-statistic as
\begin{align}
  F
  &=\frac{\operatorname{TSS}-\operatorname{RSS}}{\operatorname{RSS}/(n-2)} \\
  &=\frac{R^2\operatorname{TSS}}{(1-R^2)\operatorname{TSS}/(n-2)} \\
  &=\boxed{\frac{R^2}{1-R^2}(n-2).}
  \label{eq:ch6-f-in-terms-r2}
\end{align}

This formula explains why sample size matters.  If two simple linear regressions have the same $R^2$, the one with more observations will have a larger $F$-statistic.  More data can make the same proportion of explained variation more statistically detectable.

\section{The $F$-Distribution and the p-Value}

Under Assumption 3 from Chapter 5---the Gaussian error model---and the null hypothesis $H_0:\beta_1=0$,
\begin{equation}
  F \sim F_{1,n-2}.
\end{equation}
The first degree of freedom is the numerator degree of freedom:
\begin{equation}
  p_{\text{full}}-p_{\text{null}}=2-1=1.
\end{equation}
The second degree of freedom is the residual degree of freedom:
\begin{equation}
  n-p_{\text{full}}=n-2.
\end{equation}

After calculating the observed statistic $F_{\text{obs}}$, the p-value is
\begin{equation}
  p\text{-value}=\Pr\left(F_{1,n-2}\geq F_{\text{obs}}\mid H_0\right).
\end{equation}

A small p-value means that, if the true slope were zero and the model assumptions held, an improvement this large would be unlikely by random sampling variation alone.

\section{Why the Assumptions Matter}

The $F$-test relies on the same three-step regression assumption ladder introduced in Chapter 5.

\begin{enumerate}
  \item Assumption 1 (mean-zero errors) supports the interpretation of the line as the conditional mean.
  \item Assumption 2 (constant variance and uncorrelated errors) supports the standard error and variance calculations.
  \item Assumption 3 (Gaussian errors) supports exact finite-sample $t$ and $F$ distributions.
\end{enumerate}

If these assumptions are badly violated, the p-value may not have the advertised interpretation.  This is why residual analysis is not optional.

\section{The Slope $t$-Test}

Most software output also reports a $t$-test for the slope.  The slope test is
\begin{equation}
  H_0:\beta_1=0
  \qquad\text{versus}\qquad
  H_A:\beta_1\neq 0.
\end{equation}

The residual variance estimator is
\begin{equation}
  \widehat{\sigma}^2=\frac{\operatorname{RSS}}{n-2}.
  \label{eq:ch6-sigmahat}
\end{equation}
The standard error of the slope is
\begin{equation}
  \operatorname{SE}(\widehat{\beta}_1)
  =\sqrt{\frac{\widehat{\sigma}^2}{\sum_{i=1}^n (x_i-\bar{x})^2}}
  =\sqrt{\frac{\widehat{\sigma}^2}{S_{xx}}}.
  \label{eq:ch6-se-beta1}
\end{equation}
The $t$-statistic is
\begin{equation}
  t_{\text{obs}}=\frac{\widehat{\beta}_1}{\operatorname{SE}(\widehat{\beta}_1)}.
  \label{eq:ch6-t-stat}
\end{equation}
Under the null hypothesis and Assumption 3 (Gaussian errors),
\begin{equation}
  t_{\text{obs}}\sim t_{n-2}.
\end{equation}

\CoreCompress{
\section{Why the SLR $F$-Test and Slope $t$-Test Agree}

In simple linear regression with an intercept,
\begin{equation}
  \boxed{F=t^2.}
\end{equation}

To see why, recall that
\begin{equation}
  \operatorname{RegSS}=\widehat{\beta}_1^2S_{xx}.
\end{equation}
Also,
\begin{equation}
  \operatorname{SE}(\widehat{\beta}_1)^2=\frac{\widehat{\sigma}^2}{S_{xx}},
  \qquad
  \widehat{\sigma}^2=\frac{\operatorname{RSS}}{n-2}.
\end{equation}
Thus,
\begin{align}
  t^2
  &=\frac{\widehat{\beta}_1^2}{\operatorname{SE}(\widehat{\beta}_1)^2} \\
  &=\frac{\widehat{\beta}_1^2}{\widehat{\sigma}^2/S_{xx}} \\
  &=\frac{\widehat{\beta}_1^2S_{xx}}{\widehat{\sigma}^2} \\
  &=\frac{\operatorname{RegSS}}{\operatorname{RSS}/(n-2)} \\
  &=F.
\end{align}

Therefore, in SLR, the overall $F$-test for the model and the two-sided $t$-test for the slope produce the same conclusion.  In multiple linear regression, the overall $F$-test and individual coefficient $t$-tests answer different questions.
}{
\section{Why the SLR $F$-Test and Slope $t$-Test Agree}

In simple linear regression with an intercept,
\begin{equation}
  F=t^2.
\end{equation}
Both tests answer the same question: is the slope different from zero?  The $F$-test compares the fitted line to the mean-only model, while the slope $t$-test standardizes $\widehat{\beta}_1$ by its standard error.  In MLR these are no longer the same question: the overall $F$-test asks whether at least one slope is useful, while individual $t$-tests ask about one coefficient at a time.
}
\CoreCompress{
\section{Worked Example}

Consider five observations:
\begin{center}
\begin{tabular}{cc}
\hline
$x_i$ & $y_i$ \\
\hline
1 & 2 \\
2 & 3 \\
3 & 5 \\
4 & 4 \\
5 & 6 \\
\hline
\end{tabular}
\end{center}

The least squares line is
\begin{equation}
  \widehat{y}=1.3+0.9x.
\end{equation}
The fitted values and residuals are:
\begin{center}
\begin{tabular}{ccccc}
\hline
$i$ & $x_i$ & $y_i$ & $\widehat{y}_i$ & $e_i$ \\
\hline
1 & 1 & 2 & 2.2 & $-0.2$ \\
2 & 2 & 3 & 3.1 & $-0.1$ \\
3 & 3 & 5 & 4.0 & $1.0$ \\
4 & 4 & 4 & 4.9 & $-0.9$ \\
5 & 5 & 6 & 5.8 & $0.2$ \\
\hline
\end{tabular}
\end{center}

The sample mean is
\begin{equation}
  \bar{y}=4.
\end{equation}
The total sum of squares is
\begin{equation}
  \operatorname{TSS}=(2-4)^2+(3-4)^2+(5-4)^2+(4-4)^2+(6-4)^2=10.
\end{equation}
The residual sum of squares is
\begin{align}
  \operatorname{RSS}
  &=(-0.2)^2+(-0.1)^2+(1.0)^2+(-0.9)^2+(0.2)^2 \\
  &=1.9.
\end{align}
Therefore,
\begin{equation}
  R^2=1-\frac{1.9}{10}=0.81.
\end{equation}
The model explains $81\%$ of the observed variation in $y$.

The $F$-statistic is
\begin{align}
  F
  &=\frac{\operatorname{TSS}-\operatorname{RSS}}{\operatorname{RSS}/(n-2)} \\
  &=\frac{10-1.9}{1.9/3} \\
  &\approx 12.79.
\end{align}
Under $H_0:\beta_1=0$, this is compared to an $F_{1,3}$ distribution.  Equivalently, the slope $t$-statistic is
\begin{equation}
  t=\sqrt{F}\approx 3.58,
\end{equation}
with the sign determined by the sign of $\widehat{\beta}_1$.  Since the slope is positive, $t\approx 3.58$.
}{
\section{Worked Example}

For the five observations in the full reference example, the fitted line is
\begin{equation}
  \widehat{y}=1.3+0.9x.
\end{equation}
The total sum of squares is \(\operatorname{TSS}=10\), the residual sum of squares is \(\operatorname{RSS}=1.9\), and therefore
\begin{equation}
  R^2=1-\frac{1.9}{10}=0.81.
\end{equation}
The model explains \(81\%\) of the observed variation in \(y\).  The corresponding $F$-statistic is about \(12.79\), compared to an \(F_{1,3}\) distribution under \(H_0:\beta_1=0\).  The full reference text shows the residual table and arithmetic details.
}
\section{How to Read Software Output}

Regression software usually reports both fit-quality and hypothesis-testing information.

\subsection{Fit Quality Fields}

Common fields include:
\begin{itemize}
  \item \textbf{R-squared:} the $R^2$ value;
  \item \textbf{Residual standard error} or \textbf{scale}: an estimate of residual noise;
  \item \textbf{F-statistic:} the model-versus-null comparison;
  \item \textbf{Prob(F-statistic):} the p-value for the $F$-test.
\end{itemize}

\subsection{Coefficient Table Fields}

For the predictor row, common fields include:
\begin{itemize}
  \item \textbf{coef:} $\widehat{\beta}_1$;
  \item \textbf{std err:} $\operatorname{SE}(\widehat{\beta}_1)$;
  \item \textbf{t:} the slope $t$-statistic;
  \item \textbf{P>|t|:} the two-sided p-value for $H_0:\beta_1=0$;
  \item \textbf{confidence interval:} an interval estimate for the slope.
\end{itemize}

In SLR, the p-value in the predictor row and the p-value for the model $F$-statistic correspond to the same slope question.

\CoreCompress{
\section{Decision Template for SLR Fit and Significance}

When reporting a simple linear regression, use the following template.

\begin{enumerate}
  \item \textbf{State the fitted model:}
  \begin{equation}
    \widehat{y}=\widehat{\beta}_0+\widehat{\beta}_1x.
  \end{equation}
  \item \textbf{Interpret the slope with units.}
  \item \textbf{Report $R^2$:} say what proportion of response variation is explained by the predictor.
  \item \textbf{State the hypothesis test:}
  \begin{equation}
    H_0:\beta_1=0
    \qquad\text{versus}\qquad
    H_A:\beta_1\neq 0.
  \end{equation}
  \item \textbf{Report the test statistic and p-value.}
  \item \textbf{Translate the conclusion into context.}
  \item \textbf{Check residuals before overclaiming.}
  \item \textbf{Avoid causal language unless the design supports causation.}
\end{enumerate}
}{
\section{Decision Template for SLR Fit and Significance}

When reporting a simple linear regression, keep the distinction between fit quality and statistical evidence clear:
\begin{enumerate}
  \item state the fitted model \(\widehat{y}=\widehat{\beta}_0+\widehat{\beta}_1x\);
  \item interpret the slope with units;
  \item report \(R^2\) as an in-sample fit measure;
  \item state the slope hypothesis test \(H_0:\beta_1=0\) versus \(H_A:\beta_1\neq 0\);
  \item report the test statistic and p-value;
  \item check residuals before making final claims.
\end{enumerate}
}
\paragraph{Coding companion reference.}
Package-specific Python/R commands for computing \(R^2\), \(F\)-statistics, p-values, and related output have been moved to the separate Coding and Software Cookbook companion.  The main chapter keeps the conceptual guidance for reading software output.

\section{Common Mistakes}

\begin{enumerate}
  \item \textbf{Treating high $R^2$ as proof of causation.}  $R^2$ is about explained variation, not causal mechanism.
  \item \textbf{Treating high $R^2$ as proof of validity.}  A model can have high $R^2$ and still violate assumptions.
  \item \textbf{Ignoring sample size.}  The same $R^2$ can produce different statistical evidence depending on $n$.
  \item \textbf{Confusing residual variance with response variance.}  $\operatorname{TSS}$ measures response variation before the model; $\operatorname{RSS}$ measures variation left after the model.
  \item \textbf{Forgetting degrees of freedom.}  In SLR, residual variance uses denominator $n-2$ because two parameters are estimated.
  \item \textbf{Thinking $R^2=1$ always means a useful model.}  Two points can give a perfect line, but not a reliable model.
  \item \textbf{Reading p-values without context.}  A small p-value does not tell you effect size, practical importance, or model adequacy.
\end{enumerate}

\CoreCompress{
\section{Chapter Summary}

The mean-only model predicts every response using $\bar{y}$.  Its squared error is the total sum of squares,
\begin{equation}
  \operatorname{TSS}=\sum_{i=1}^n (y_i-\bar{y})^2.
\end{equation}
The fitted regression model predicts each response with
\begin{equation}
  \widehat{y}_i=\widehat{\beta}_0+\widehat{\beta}_1x_i.
\end{equation}
Its squared error is the residual sum of squares,
\begin{equation}
  \operatorname{RSS}=\sum_{i=1}^n (y_i-\widehat{y}_i)^2.
\end{equation}
The coefficient of determination is
\begin{equation}
  R^2=1-\frac{\operatorname{RSS}}{\operatorname{TSS}}.
\end{equation}
It measures the proportion of observed response variation explained by the fitted line.

The $F$-test compares the fitted line to the mean-only model:
\begin{equation}
  F=\frac{\operatorname{TSS}-\operatorname{RSS}}{\operatorname{RSS}/(n-2)}.
\end{equation}
Under $H_0:\beta_1=0$ and Assumption 3 (Gaussian errors),
\begin{equation}
  F\sim F_{1,n-2}.
\end{equation}
In simple linear regression with an intercept, the slope $t$-test and model $F$-test are equivalent:
\begin{equation}
  F=t^2.
\end{equation}

The central distinction is this:
\begin{quote}
  $R^2$ measures how much the model fits the observed data; the $F$-test and $t$-test measure whether the observed improvement is statistically distinguishable from a zero-slope model under the same regression assumption ladder.
\end{quote}

\section{Math Foundations Source Map}

This source map records the Math Foundations support used in this chapter and where those ideas appear in the completed Math Foundations textbook.

\begin{center}
\begin{tabular}{|p{0.24\textwidth}|p{0.36\textwidth}|p{0.28\textwidth}|}
\hline
\textbf{Chapter 6 use} & \textbf{Math Foundations connection} & \textbf{MF textbook location} \\
\hline
Mean-only fitted vector & $\mathbb{R}^n$ vector notation and the ones vector $\mathbf{1}$ & Ch02 Secs. 2.1, 2.4 \\
\hline
TSS and RSS as squared lengths & Dot products, $\mathbf{a}^T\mathbf{a}=\sum_i a_i^2$, and Euclidean norm $\|\mathbf{x}\|=\sqrt{\mathbf{x}^T\mathbf{x}}$ & Ch05 Secs. 5.1, 5.3, 5.5 \\
\hline
Sum-of-squares decomposition & Dot products, norms, and orthogonality of fitted and residual components & Ch02 Sec. 2.6; Ch05 Secs. 5.1, 5.3, 5.5, 5.8; Ch12 Sec. 12.12 \\
\hline
Residual degrees of freedom & Linear independence, basis size, and leftover dimension intuition & Ch02 Sec. 2.6; Ch06 Secs. 6.3, 6.4, 6.5 \\
\hline
Hypothesis statements & Logic statements, negation, and truth-table reasoning as the foundation for null-versus-alternative statements & Ch01 Sec. 1.1 \\
\hline
Software/design-matrix bridge & Matrix notation, transpose notation, and matrix-vector multiplication & Ch07 Secs. 7.1, 7.4, 7.6.1, 7.6.2 \\
\hline
\end{tabular}
\end{center}

\section{Practice and Workflow}
\subsection*{Focused Study Checklist}

Use this as a final check after analyzing fit quality and statistical evidence.

\begin{enumerate}
  \item Can you identify the mean-only baseline and compute \(\operatorname{TSS}\)?
  \item Can you compute residuals, \(\operatorname{RSS}\), and \(R^2=1-\operatorname{RSS}/\operatorname{TSS}\)?
  \item Can you explain \(R^2\) as fit quality, not as proof that the model is useful or causal?
  \item Can you state the slope test hypotheses \(H_0:\beta_1=0\) and \(H_A:\beta_1\neq 0\)?
  \item Can you connect the model \(F\)-test to the slope \(t\)-test in SLR through \(F=t^2\)?
  \item Can you conclude in context while separating fit, statistical evidence, and residual diagnostics?
\end{enumerate}
}{
\section{Chapter Summary}

The mean-only model predicts every response by \(\bar{y}\).  Its squared error is \(\operatorname{TSS}\).  The regression model predicts each response by \(\widehat{y}_i=\widehat{\beta}_0+\widehat{\beta}_1x_i\).  Its remaining squared error is \(\operatorname{RSS}\).  The key fit measure is
\begin{equation}
  R^2=1-\frac{\operatorname{RSS}}{\operatorname{TSS}},
\end{equation}
which is the proportion of observed response variation explained by the fitted line.

The $F$-test compares the fitted line to the mean-only model.  Under \(H_0:\beta_1=0\) and the Gaussian-error rung of the regression assumption ladder,
\begin{equation}
  F=\frac{\operatorname{TSS}-\operatorname{RSS}}{\operatorname{RSS}/(n-2)}\sim F_{1,n-2}.
\end{equation}
In SLR with an intercept, the slope $t$-test and model $F$-test agree because \(F=t^2\).

\paragraph{Can you do this?}
You should be able to compute TSS, RSS, \(R^2\), and the SLR $F$-statistic; explain the difference between fit and statistical evidence; and read the slope row and model $F$-test from software output.  The full reference text contains the complete Math Foundations source map and reusable study template.
}
